\documentclass[11pt,a4paper]{article}

\usepackage{fontspec}
\usepackage{polyglossia}
\setdefaultlanguage{polish}

\usepackage[a4paper,margin=2.5cm]{geometry}
\usepackage[backend=biber,style=numeric,sorting=none]{biblatex}
\usepackage{graphicx}
\usepackage{hyperref}
\usepackage{booktabs}
\usepackage{amsmath}
\usepackage{enumitem}

\addbibresource{bibliography.bib}

\hypersetup{
    colorlinks=true,
    linkcolor=black,
    citecolor=blue,
    urlcolor=blue
}

\title{Zastosowanie MCTS/UCT do gry Breakthrough 6$\times$6 \\
       \large Konspekt projektu --- MSI, semestr letni 2025/26}
\author{Sebastian Rydzewski, Hubert [Nazwisko]}
\date{\today}

\begin{document}
\maketitle

\section{Opis problemu}

\paragraph{Gra Breakthrough.}
Breakthrough to dwuosobowa gra strategiczna z pełną informacją, wynaleziona
przez Dana Trokę w 2000 roku (zwycięzca konkursu 8x8 Game Design Competition).
Plansza jest prostokątna; każdy gracz dysponuje pionkami ustawionymi w dwóch
pierwszych rzędach po swojej stronie. Ruch pionka odbywa się o jedno pole do
przodu --- prosto lub po skosie. Ruch prosty jest możliwy wyłącznie na puste
pole; ruch po skosie może być wykonany na pole puste lub zajęte przez pionek
przeciwnika (bicie). Bicie prosto jest niedozwolone, a pionki nie mogą się
cofać. Gra kończy się zwycięstwem gracza, który: (a) doprowadzi pionek do
ostatniego rzędu przeciwnika, (b) zbije wszystkie pionki przeciwnika, lub
(c) uniemożliwi przeciwnikowi wykonanie jakiegokolwiek legalnego ruchu.
W klasycznych regułach nie występują remisy.

\paragraph{Dlaczego Breakthrough jako benchmark MCTS.}
Gra spełnia wszystkie wymagania projektu:
\begin{itemize}[nosep]
    \item \textbf{Dwuosobowa, z pełną informacją, deterministyczna} ---
          kluczowe właściwości dla metod MCTS/UCT bez rozszerzeń probabilistycznych.
    \item \textbf{Umiarkowany branching factor} --- w typowej pozycji liczba
          legalnych ruchów wynosi 10--20. Głębokość partii na planszy 6$\times$6
          to 20--40 półruchów.
    \item \textbf{Niebanalna strategicznie} --- mimo prostych zasad wymaga planowania
          na wiele ruchów naprzód (tworzenie łańcuchów pionków, kontrola kolumn,
          poświęcenia dla osiągnięcia przełomu).
    \item \textbf{Niszowa poza środowiskiem badawczym} --- w odróżnieniu od
          szachów czy Go nie istnieją silne silniki oparte o sieci neuronowe,
          co czyni ją atrakcyjnym polem dla porównań klasycznych wariantów MCTS.
    \item \textbf{Wariant 6$\times$6 nie został rozwiązany} --- Lorentz rozwiązał
          warianty do planszy 5$\times$5 \cite{saffidine2011solving}, ale 6$\times$6
          pozostaje otwarty, co zostawia miejsce na eksperymenty porównawcze bez
          znanego ,,perfekcyjnego gracza'' ustawiającego sufit jakości.
\end{itemize}

\paragraph{Modyfikacja reguł: plansza 6$\times$6.}
Zdecydowaliśmy się na wariant 6$\times$6 zamiast klasycznej 8$\times$8 z trzech powodów.
Po pierwsze, mniejsza plansza pozwala na większą liczbę iteracji MCTS w tym samym
budżecie czasowym --- co jest istotne przy wymogu minimum 30 partii na konfigurację
i 4 hipotezach do zweryfikowania. Po drugie, 6$\times$6 zachowuje charakter strategiczny
gry (branching factor $\sim$10--20 pozostaje w przedziale przyjaznym MCTS), jednocześnie
umożliwiając reprezentację planszy w~pojedynczym bitboardzie \texttt{u64} (36 pól),
co upraszcza implementację generatora ruchów. Po trzecie, 6$\times$6 nie jest rozwiązany,
co oznacza że porównanie wariantów MCTS ma realną wartość badawczą.

\section{Wstępny przegląd literatury}

Metoda Monte Carlo Tree Search została zaproponowana przez
Coulom~\cite{coulom2006efficient} jako algorytm przeszukiwania dla gry Go.
W tym samym roku Kocsis i Szepesvári~\cite{kocsis2006bandit} opublikowali
algorytm UCT (Upper Confidence Bounds applied to Trees), łączący klasyczne
MCTS z regułą wyboru UCB1 znaną z teorii bandytów. Kompleksowy przegląd
metody i jej wariantów przedstawili Browne et al.~\cite{browne2012survey}
--- jest to główny punkt odniesienia metodologicznego dla niniejszego projektu.

\paragraph{RAVE i AMAF.}
Gelly i Silver~\cite{gelly2011monte} wprowadzili algorytm RAVE (Rapid Action
Value Estimation), rozszerzający UCT o dodatkowe statystyki AMAF (All-Moves-As-First).
Kluczowa intuicja: jeśli ruch $m$ okazywał się dobry w symulacjach wychodzących
z potomków węzła $n$, to $m$ jest również prawdopodobnie dobry z samego $n$.
RAVE wyraźnie przyspiesza zbieżność UCT w grach, w~których te same ruchy pojawiają
się w różnych kontekstach --- co ma miejsce w Breakthrough (pionek na kolumnie $c$
może zostać zagrany na tym samym polu w różnych partiach drzewa).

\paragraph{Progressive Bias.}
Chaslot et al.~\cite{chaslot2008progressive} zaproponowali rodzinę technik
``Progressive Strategies'', z których jedna --- Progressive Bias --- wstrzykuje
domenową heurystykę do fazy selekcji MCTS. Wartość heurystyki
$H(m)$ dodawana jest do klasycznej formuły UCB1 z wagą malejącą wraz z~liczbą
odwiedzin węzła: $H(m)/(n+1)$. Technika łączy silne strony czystego MCTS
(zbieżność do optimum) z szybkim startem heurystyki (dobre ruchy w pierwszych
iteracjach).

\paragraph{Heurystyki specjalistyczne dla Breakthrough.}
Lorentz~\cite{lorentz2011breakthrough} opisał techniki ewaluacji pozycji dla
Breakthrough, w szczególności wartościowanie zaawansowania pionków, bezpieczeństwa
na krawędziach oraz mobilności. Jego praca ,,Early Playout Termination in MCTS''
\cite{lorentz2016early} pokazuje że rozsądna heurystyka potrafi być konkurencyjna
z MCTS o niskim budżecie iteracji.

\paragraph{Pierwszy gracz a rozmiar planszy.}
Saffidine et al.~\cite{saffidine2011solving} rozwiązali Breakthrough dla planszy
5$\times$5 i~wykazali że białe (pierwszy gracz) wygrywają przy perfekcyjnej grze.
Hipotetyczna przewaga pierwszego gracza w Breakthrough jest przedmiotem dyskusji
w społeczności AI, ale brakuje systematycznych eksperymentów empirycznych dla
średnich wariantów (6$\times$6, 6$\times$8, 8$\times$8).

Wszystkie cytowane publikacje są recenzowane (konferencje AIIDE, Computer Games,
IEEE Transactions on Games). Zgodnie z wymaganiem projektu unikamy blogów
i niezrecenzowanych preprintów.

\section{Hipotezy badawcze}

Zgodnie z wymaganiami projektu stawiamy cztery hipotezy. Pierwsza dotyczy
rozszerzenia UCT z literatury, druga --- drugiego wariantu metodologicznie
różnego od pierwszego, trzecia --- specjalistycznej heurystyki, czwarta ---
własnej hipotezy związanej z zasadami gry.

\paragraph{H1 --- RAVE poprawia jakość gry względem vanilla UCT.}
Hipotezujemy, że UCT wzbogacone o RAVE osiąga istotnie statystycznie wyższy
odsetek wygranych niż vanilla UCT przy tym samym budżecie iteracji.
\emph{Oczekiwany kierunek:} poprawa rzędu 5--15 punktów procentowych przy
budżecie 1000--10000 iteracji, malejąca z~budżetem (przy bardzo wysokim
budżecie UCT samo zbiega do dobrego rozwiązania).

\paragraph{H2 --- Progressive Bias poprawia jakość gry względem vanilla UCT.}
Hipotezujemy, że UCT z Progressive Bias (wykorzystujące tę samą heurystykę
co gracz H3) osiąga istotnie wyższy odsetek wygranych niż vanilla UCT przy
tym samym budżecie iteracji. Wybraliśmy Progressive Bias zamiast alternatyw
(MAST, decisive moves) ze względu na: (a) solidne oparcie w literaturze,
(b) metodologiczne rozróżnienie od RAVE --- PB wprowadza wiedzę domenową
do selekcji, RAVE wyłącznie statystyki ze symulacji, (c) naturalne
powiązanie z H3: ta sama funkcja heurystyczna zostanie użyta w dwóch
eksperymentach, co pozwoli na spójną analizę.

\paragraph{H3 --- Specjalistyczna heurystyka pokonuje MCTS przy niskim budżecie.}
Hipotezujemy, że gracz oparty na minimaxie alfa-beta głębokości 5 z funkcją
ewaluacji łączącą zaawansowanie pionków, materiał, zagrożenia i mobilność,
wygrywa z~vanilla UCT przy niskich budżetach iteracji ($\leq$ 5000), ale przegrywa
przy wysokich budżetach ($\geq$ 50000). Oczekujemy zatem istnienia ,,punktu krzyżowania''
(cross-over) na krzywej win-rate vs budżet iteracji. Wynik potwierdzający pokaże
klasyczny trade-off: wiedza domenowa jest wartościowa przy ograniczonych zasobach,
ale MCTS dominuje przy wystarczającej mocy obliczeniowej.

\paragraph{H4 --- Przewaga pierwszego gracza maleje ze wzrostem rozmiaru planszy.}
Hipotezujemy, że białe (pierwszy gracz) mają istotnie statystyczną przewagę
w Breakthrough 6$\times$6 przy równej sile agentów MCTS, i że ta przewaga maleje
w kierunku 8$\times$8. Konkretnie: dla turnieju UCT 5000 iter vs UCT 5000 iter,
win rate białych będzie wyraźnie powyżej 50\% dla 6$\times$6, zbliżać się do 50\%
dla 6$\times$8, i być jeszcze bliższy 50\% dla 8$\times$8. Hipoteza jest motywowana
obserwacją że na mniejszej planszy pierwszy ruch daje bezpośrednią przewagę
tempa, podczas gdy większa plansza rozcieńcza efekt pierwszego ruchu.

\section{Sposób weryfikacji hipotez}

\paragraph{Eksperymenty.}
Każda hipoteza będzie testowana w oddzielnej serii turniejów:
\begin{itemize}[nosep]
    \item \textbf{H1:} round-robin RAVE vs UCT dla budżetów
          $\{1000, 5000, 10000, 50000\}$ iteracji.
    \item \textbf{H2:} round-robin Progressive Bias vs UCT dla tych samych budżetów.
    \item \textbf{H3:} round-robin gracz heurystyczny (głębokość 5) vs UCT
          dla budżetów $\{200, 500, 1000, 5000, 10000, 50000\}$.
    \item \textbf{H4:} UCT 5000 iter vs UCT 5000 iter dla trzech rozmiarów planszy:
          6$\times$6, 6$\times$8, 8$\times$8.
\end{itemize}

\paragraph{Metodologia statystyczna.}
Dla każdej pary konfiguracji rozegramy minimum 60 partii (30 z~każdą stroną jako
białe). W H4 zwiększamy do 100 partii na konfigurację ze względu na skupienie
hipotezy na przewadze tempa. Z~wyników wyliczamy:
\begin{itemize}[nosep]
    \item Win rate z 95\% przedziałem ufności Wilsona (preferowany nad Waldem
          dla proporcji bliskich ekstremom).
    \item Rozkład długości partii (mediana, IQR, boxplot).
    \item Średni czas na ruch dla każdego agenta.
\end{itemize}
Uznamy hipotezę za potwierdzoną statystycznie jeśli 95\% CI win rate nie przecina
50\% w spodziewanym kierunku dla co najmniej trzech testowanych budżetów.

\paragraph{Reprodukowalność.}
Każde uruchomienie eksperymentu ma jawny master seed podawany z CLI.
Z master seed deterministycznie wyprowadzamy seedy indywidualnych partii
($\texttt{seed}_i = \text{hash}(\text{master}, \text{game\_id})$), dzięki czemu
każda partia jest niezależnie reprodukowalna. Każda partia loguje swój seed,
konfigurację, wynik, długość i czasy do~pliku JSONL (format crash-safe, append-only,
umożliwia wznowienie przerwanego eksperymentu).

\paragraph{Opracowanie wyników.}
Zamiast surowych tabel z pojedynczymi wartościami raportujemy:
\begin{itemize}[nosep]
    \item Krzywe skalowania win-rate vs budżet iteracji (z zaznaczonymi przedziałami ufności)
          --- dla H1, H2, H3.
    \item Boxploty długości partii per konfiguracja.
    \item Heatmapy win-rate dla H4 (rozmiar planszy $\times$ strona).
\end{itemize}
Wszystkie wykresy generowane będą skryptami Pythona (matplotlib + seaborn)
eksportującymi do PDF do katalogu \texttt{report/figures/}.

\paragraph{Testy z udziałem osób.}
Zgodnie z wymaganiem projektu, interfejs graficzny (Pygame) zostanie przetestowany
przez minimum 5 osób. Każda partia człowiek-komputer zostanie zalogowana
(identyfikator testera, konfiguracja agenta, wszystkie ruchy, wynik).
Raportowany wynik: liczba partii, rozkład wyników, subiektywna ocena siły agenta.

\section{Harmonogram}

Projekt realizowany od kwietnia do czerwca 2026. Zakładamy 9 tygodni pracy
z dwutygodniowym buforem na nieprzewidziane problemy.

\begin{center}
\begin{tabular}{p{4cm}p{10cm}}
\toprule
\textbf{Tydzień} & \textbf{Zadanie} \\
\midrule
Tydz. 1 & Konspekt (ten dokument), akceptacja prowadzącego \\
Tydz. 2 & Silnik gry w Ruście + referencyjny w Pythonie + testy cross-validation \\
Tydz. 3 & Vanilla UCT + testy \\
Tydz. 4 & RAVE + Progressive Bias \\
Tydz. 5 & Heurystyczny gracz (alfa-beta głębokość 5) + harness eksperymentów \\
Tydz. 6 & Uruchomienie eksperymentów H1--H4 (w tle) \\
Tydz. 7 & GUI + testy z użytkownikami \\
Tydz. 8 & Analiza wyników + generowanie wykresów \\
Tydz. 9 & Pisanie raportu + prezentacja \\
\bottomrule
\end{tabular}
\end{center}

\section{Ogólny projekt techniczny}

\paragraph{Stos technologiczny.}
Rdzeń obliczeniowy --- silnik gry, wszystkie warianty MCTS, heurystyka ---
implementujemy w języku Rust. Orchestracja eksperymentów, interfejs graficzny
i analiza wyników --- w języku Python 3.12. Wiązanie Rust$\leftrightarrow$Python
poprzez PyO3 z narzędziem maturin (build). Wybór Rusta motywowany wydajnością
(MCTS z dużym budżetem iteracji dominuje w czasie obliczeń) i bezpieczeństwem
pamięci.

\paragraph{Reprezentacja planszy.}
Dwa bitboardy \texttt{u64} (białe, czarne). Bit $i$ odpowiada polu $i$ przy
numeracji wiersz-kolumna. Generator ruchów zaimplementowany operacjami bitowymi
(przesunięcia $\pm$cols, maski krawędzi). Stan gry jest niemutowalny ---
\texttt{apply(mv)} zwraca nową instancję \texttt{GameState} (ok. 32 bajtów).
Upraszcza to kod MCTS kosztem drobnych alokacji (akceptowalny narzut przy
rozmiarze struktury i puli 50\,000 iteracji).

\paragraph{Zrównoleglenie.}
Kluczowa decyzja metodologiczna: \textbf{MCTS działa sekwencyjnie, zrównoleglenie
jest wyłącznie na poziomie batchowania eksperymentów}. Każda partia turnieju to
osobny proces Python (\texttt{ProcessPoolExecutor}). Nie wprowadzamy tree
parallelization ani root parallelization. Uzasadnienie: czystość metodologiczna
porównań --- porównujemy warianty algorytmiczne, nie wydajność zrównoleglenia.
Budżet MCTS mierzony liczbą iteracji pozostaje jednoznaczny.

\paragraph{Reprodukowalność.}
Generator liczb losowych: \texttt{Xoshiro256PlusPlus} (seedowalny), po jednym
niezależnym strumieniu na agenta. Seed agenta wyprowadzany z master seed plus
identyfikatora partii. Wszystkie hiperparametry (stała eksploracji $c$, $k$ w RAVE,
waga bias w PB, głębokość alfa-bety, wagi ewaluacji) konfigurowalne
z linii poleceń, nie hardkodowane.

\paragraph{Struktura repozytorium.}
\begin{verbatim}
breakthrough-mcts/
  Cargo.toml, pyproject.toml       # build
  src/                              # Rust
    game.rs                         # bitboardy, ruchy, terminale
    mcts/{uct,rave,progressive_bias}.rs
    heuristic.rs                    # alfa-beta
  python/breakthrough/              # Python
    reference_game.py               # czytelna implementacja do walidacji
    gui.py                          # Pygame
    experiments/{harness,run_h1..h4}.py
  tests/test_rust_vs_python.py      # cross-validation (wymagany gate)
  notebooks/                        # skrypty analityczne
  report/                           # raport LaTeX
  konspekt/                         # ten dokument
  results/                          # JSONL (w .gitignore)
\end{verbatim}

\paragraph{Walidacja poprawności silnika.}
Przed rozpoczęciem eksperymentów uruchamiamy test cross-validation:
1000 losowych partii rozegranych równolegle przez silnik Rust i Python ---
muszą generować identyczne zbiory legalnych ruchów, stany terminalne i zwycięzców
w każdym kroku. Niepoprawny silnik unieważnia wszystkie wyniki eksperymentalne.

\paragraph{Użycie LLM.}
Kod generowany przy współpracy z Claude Opus 4.7 --- każdy wygenerowany fragment
podlega review autorów. W raporcie końcowym zamieścimy szczegółową sekcję
dokumentującą sposób użycia LLM (konkretne promptty, wykryte błędy w wygenerowanym
kodzie, poprawki wprowadzone manualnie).

\printbibliography

\end{document}
